\scanpage{179}

\deleted{On peut supposer $X$ irréductible de pt générique $\eta$.
L'inclusion $\unsure{\eta-X|\eta\to X}$ définit une suite spectrale
(pour la cohomologie étale)
\[
H^*(\eta,G_m)\Leftarrow E_2^{pq}=H^p(X,R^qi_*(G_m))
\]
Or \deleted{\unsure{comme les}} comme les $H^q(\eta,G_m)$ et les fibres de $R^qi_*(G_m)$, qui s'explicitent par de la cohomologie galoisienne, sont des groupes de torsion pour $q>0$,
\unsure{je trouve que} \deleted{\unsure{il en va de même des} \unsure{gpes|groupes}}
\[
(*)\qquad H^p(X,R_X^*)\quad
\deleted{,\quad R_X^*=R\unsure{i_*}(G_m)}
\]
\inserted{est un groupe de torsion pour $p>0$,}
où $R_X^*=\deleted{\unsure{Ri_*}}i_*(G_m)$ est le faisceau (étale) des fonctions rationnelles \inserted{non nulles} sur $X$. \deleted{(\unsure{ceci étant}}}

D'autre part, on a \inserted{si $X$ est \inserted{noethérien} ``sans idéaux premiers immergés'',} la suite exacte
\diagram{\[
\text{(1.1.)}\qquad
\begin{tikzcd}[column sep=large]
0 \arrow[r,dash] & G_{mX} \arrow[r] & R_X^* \arrow[r] & \underline{\operatorname{Div}}_X \arrow[r] & 0
\end{tikzcd}
\]}
où $\underline{\operatorname{Div}}_X$ est le faisceau (\inserted{pour la top.} étale) des diviseurs (de Cartier) sur $X$, ce qui \deleted{\unsure{veut dire que}} \inserted{implique} compte tenu de ce que
\diagram{\[
\begin{tikzcd}[column sep=large]
H^{\unsure{i}}(X,\underline{\operatorname{Div}}_X) \arrow[r,"\partial"] & H^{i+1}(X,G_{mX})
\end{tikzcd}
\]}
\margin{\uline{\unsure{donne le|donner le}} \uline{corollaire} \unsure{3.4.}}
est bijectif modulo torsion \deleted{pour} si $i\geq 1$. \uline{Ainsi}
\deleted{\unsure{pour en montrer l'utilité}} L'hypothèse de factorialité d'autre part implique \deleted{l'isomorphisme} qu'on a un isomorphisme de faisceaux étales :

\margin{178}
